\section{Conclusion}
\label{sec:ee:conclusion}

This chapter addressed \ref{ob:extraction_assessment} by empirically assessing the applicability of general-domain \ac{ee} methods to legal and investigative contexts. In the first part (\Cref{sec:ijckg}), we introduced the incremental \ac{nel} problem, developed a methodology to adapt static \ac{nel} datasets to the incremental setting, and released both a benchmark and a baseline pipeline. The evaluation showed that incremental settings suffer from error propagation, and that NIL prediction is major source of errors.


In the second part, we evaluated \ac{ner} and \ac{inel} pipelines on Italian civil judgments and studied strategies for adapting \ac{ner} to the domain. Without in-domain fine-tuning, models trained on public corpora degraded on all categories except \texttt{Person}, and errors at this stage propagated to subsequent components. This confirms the need for domain-specific adaptation and motivates the \acl{hitl} strategy adopted to improve the quality of machine annotations.

Considering fine-tuning, the combination of in-domain \acl{mlm} and task-specific \ac{ner} fine-tuning achieved the best results, but simply fine-tuning a general-domain model already proved to be a good compromise between performance and computational cost. These results indicate that the size of our labeled dataset is sufficient to fine-tune effective \ac{ner} models. A natural direction for future work is to quantify the amount of labeled text needed to reach satisfactory performance for this domain.

Accurate NIL prediction remains a key challenge in \acl{ee} in the legal domain, where relevant actors are not present in public \acp{kr}, and misclassifying them as existing entities contaminating downstream consolidation. In-domain fine-tuning may mitigate this effect and represents another avenue for future work.

We subsequently extended the analysis to investigative \ac{ima} data, creating an annotated benchmark and evaluating an adapted \ac{ee} pipeline. Results indicate that entity extraction in \ac{ima} is challenging and could benefit from in-domain fine-tuning, whereas metadata can be reliably extracted to enable faceted search and other \ac{ia} functionalities.

These findings motivate, in the subsequent chapters, the design of architectures that remain useful despite imperfect extraction, thanks to traceability and error-correction capability.
